\subsection{Otimização de Memória (\texttt{--mem-opt})}
\label{sec:memoria}

Na formulação ingênua com matriz completa (\texttt{--full-matrix}), o consumo de memória escala quadraticamente ($\mathcal{O}(N \cdot M)$). Para matrizes de grande porte ($100.000^2$ a $320.000^2$), essa alocação exigiria de 40 GiB a centenas de gigabytes de RAM, além de transformar a execução em um problema estritamente limitado pela largura de banda da memória (\textit{memory-bound}).

Como o objetivo é obter apenas o \textit{score} escalar final (e não o alinhamento completo), as células internas podem ser descartadas à medida que o cálculo diagonal em blocos avança:
\begin{itemize}
    \item \textbf{Sequencial:} Mantém apenas dois vetores de linha alternados (\texttt{prev\_row} e \texttt{curr\_row}), reduzindo o consumo para $\mathcal{O}(N)$.
    \item \textbf{OpenMP:} Armazena apenas os vetores de fronteira (\texttt{horiz\_borders} e \texttt{vert\_borders}) entre blocos. Cada tarefa aloca um buffer local de $513 \times 513$ elementos na pilha/cache, computa o miolo e grava de volta apenas a última linha e coluna.
    \item \textbf{MPI:} Cada processo mantém apenas dois vetores de tamanho $\texttt{CHUNK\_SIZE} + 1$ e um vetor vertical \texttt{left\_border} para a fronteira entre \textit{chunks}. Isso induz um \textbf{\textit{cache blocking}} ideal: duas linhas de 1024 posições ocupam 2 a 4 KiB, residindo integralmente na cache L1d (32 KiB).
\end{itemize}
Essa otimização viabiliza matrizes gigantescas (como Q7 de $160.000^2$ e Q8 de $320.000^2$) e elimina a contenção no barramento de memória.
